\documentclass[11pt,a4paper]{article}
\usepackage[utf8]{inputenc}
\usepackage[T1]{fontenc}
\usepackage[italian]{babel}
\usepackage{graphicx}
\usepackage{booktabs}
\usepackage{amsmath}
% \usepackage{siunitx}  % non installato (texlive-mathscience); \, e \text{dB} funzionano senza
\usepackage[margin=2.5cm]{geometry}
\usepackage{hyperref}

% PRIMARIO: banda piena (48 kHz) + file intero, in analisi/distanza-banda-piena/.
% Il 10 kHz (analisi/distanza/) e la finestra 2 s restano solo come confronto.
\graphicspath{%
  {../analisi/distanza-banda-piena/livelli/}%
  {../analisi/distanza-banda-piena/centroide/}%
  {../analisi/distanza-banda-piena/descrittori/}%
  {../analisi/distanza-banda-piena/descrittori/noise/}%
  {img/report/}}

\title{Descrittori audio alla prova della distanza\\
\large Un esperimento al LEAP: stessa sorgente, otto punti di ripresa}
\author{}
\date{Roma, 13 giugno 2026}

\begin{document}
\maketitle

\begin{abstract}
Documento di laboratorio. Ho allestito al LEAP una condizione di test per misurare
quanto la risposta dei descrittori dipenda dal punto di ripresa nello spazio acustico,
con otto microfoni a distanza crescente (da 1 a 8 metri) dalla stessa sorgente. Il
centroide spettrale, invariante di scala, fa da sensore principale: se cambia tra un
microfono e l'altro non è perché il segnale è più debole, ma perché la catena
fatta di sorgente, ambiente e microfono ne ha deformato la forma. Riporto la procedura, la
calibrazione, e le discrepanze rilevate sul rumore e sugli strumenti.
\end{abstract}

\section{Scopo}
Voglio capire quanto la risposta dei descrittori dipenda dal punto in cui prendo
il segnale. Nell'uso in tempo reale do per scontato che, fissata la sorgente, i
descrittori restituiscano lo stesso vettore; ma fra la sorgente e il microfono c'è
un luogo, con il suo carattere, e una catena di trasduzione. La prova che metto in
atto è semplice: registro la stessa sorgente con otto microfoni a distanza crescente,
da 1 a 8 metri, e guardo se la descrizione resta ferma oppure si muove.

Il sensore principale è il centroide spettrale. La sua proprietà utile è che è
invariante di scala: moltiplicare lo spettro per un guadagno non lo sposta (lo mostro
nelle note matematiche in appendice). Questo
lo rende un controllo pulito: se il segnale a 8 metri è più debole che a 1 metro, il
centroide non se ne accorge. Se invece il centroide cambia da un microfono all'altro,
non è perché il livello è calato, ma perché la catena sorgente, ambiente e microfono
ha rimodellato la forma dello spettro. Trovare variazioni significherebbe che il punto
di ripresa, e non solo la sorgente, entra a pieno titolo nel valore misurato; trovarlo
piatto vorrebbe dire che il centroide è robusto rispetto al luogo.

\section{Apparato e metodo}
Ho lavorato al LEAP (Laboratorio ElettroAcustico Permanente, Roma). Per la ripresa ho
usato otto microfoni omnidirezionali LineAudio OM1. Per la calibrazione li ho prima
raggruppati con un elastico, capsule equidistanti dalla sorgente, in un unico grappolo:
in questa posizione ho regolato a mano i gain finché ogni microfono restituiva la stessa
ampiezza RMS, così da partire da una base comune (figura~\ref{fig:apparato-calibrazione}).

\begin{figure}[ht]
  \centering
  \begin{minipage}[b]{0.40\linewidth}\centering
    \includegraphics[width=\linewidth]{grappolo.jpg}
  \end{minipage}\hfill
  \begin{minipage}[b]{0.57\linewidth}\centering
    \includegraphics[width=\linewidth]{calibrazione.jpg}
  \end{minipage}
  \caption{A sinistra il grappolo degli otto microfoni omnidirezionali LineAudio OM1,
  tenuti insieme con un elastico e capsule affiancate. A destra la geometria di
  calibrazione: il grappolo posto a 1 m di fronte al diffusore Yamaha HS8, all'altezza
  mediana fra cono e tweeter, dove ho allineato i gain.}
  \label{fig:apparato-calibrazione}
\end{figure}

La sorgente di riferimento era un diffusore Yamaha HS8 su stativo, con i microfoni posti
all'altezza mediana fra cono e tweeter. La preamplificazione era affidata a un Audient
ASP880, impedenza su MED, filtri disattivati e gain regolati a mano. Al fonometro ho
letto 96 dBA (curva A) a 1 metro; ho registrato a 96 kHz, otto canali simultanei.

Per la prova vera e propria ho disposto i microfoni in fila a distanza crescente, uno
per metro: il canale $n$ a $n$ metri (ch1 a 1 m fino a ch8 a 8 m), passo costante di un
metro. Per ogni presa ho registrato due orientamenti della capsula: orizzontale, con la
capsula rivolta verso la sorgente, e verticale, con la capsula rivolta verso l'alto,
ortogonale alla direzione del suono. Lo scopo della doppia ripresa è separare due cause
che altrimenti si confonderebbero: la fisica della propagazione, che è la stessa nei due
orientamenti, e la direttività dell'OM1, che ai medio-acuti smette di comportarsi da
omnidirezionale e sposta la sensibilità verso il fronte. Se un effetto compare in
entrambi gli orientamenti non può essere della direttività del microfono; se cambia
ruotando le capsule, lo è.

Le sorgenti sono state tre. Il rumore rosa dal diffusore, come segnale di prova
controllato e ripetibile. Il clarinetto contrabbasso (Do grave a tubo chiuso e cluster),
suonato sia con il trombino originale svasato sia con un trombino custom verticale che
ho chiamato PETALONIO. Infine il clarinetto soprano (cluster)
(figura~\ref{fig:apparato-sessione}).

\begin{figure}[ht]
  \centering
  \begin{minipage}[b]{0.47\linewidth}\centering
    \includegraphics[width=\linewidth]{fonometro.jpg}
  \end{minipage}\hfill
  \begin{minipage}[b]{0.50\linewidth}\centering
    \includegraphics[width=\linewidth]{petalonio.jpg}
  \end{minipage}
  \caption{A sinistra il fonometro durante la lettura del livello a 1 m (96 dBA, curva A).
  A destra la sessione con il clarinetto contrabbasso: il trombino verticale PETALONIO
  rivolto verso l'alto e la fila delle aste microfoniche che si allontana nello spazio
  acustico.}
  \label{fig:apparato-sessione}
\end{figure}

\section{Pipeline d'analisi}
Ho analizzato i segnali con \texttt{scripts/analisi.py}. La trasformata usa una finestra
di Hann lunga 8192 campioni con passo 4096 (sovrapposizione del 50\%); a 96 kHz questo
corrisponde a frame di circa 85 ms e a una risoluzione in frequenza di circa 11,7 Hz per
bin. Analizzo l'intera banda fino a 48 kHz (la frequenza di Nyquist a 96 kHz). Un tetto a
10 kHz, usato in una prima versione, lo tengo ora solo come termine di confronto: serve a
mostrare cosa si perde tagliando gli acuti (vedi §\ref{sec:rumore}).

Per escludere il rumore di fondo dello luogo ho applicato due gate: uno assoluto a
$-65$ dBFS e uno relativo a $-30$ dB rispetto al frame più forte. Ho analizzato ogni
microfono separatamente, senza mediarli (opzione \texttt{-{}-per-canale}), perché è
proprio la differenza fra un canale e l'altro che mi interessa leggere. Su ciascun campione
medio i descrittori su tutta la parte attiva del segnale (i frame che superano i gate, non il
silenzio): uso il file intero, non una finestra corta. In una prima versione avevo usato una
finestra di 2 s; la tengo come confronto, ma i numeri qui sono sull'intero.

Da ogni frame estraggo i 16 descrittori, raggruppati in quattro famiglie: forma
(centroide, spread, rolloff, slope, obsir-std), distribuzione (flatness, crest, skewness,
kurtosis, entropy), tonalità (tpr, numero di picchi, tonality) e dinamica (flux,
irregularity, zcr). Gli output (CSV e figure) sono rigenerabili in
\texttt{analisi/distanza/}, in sottocartelle per tema (centroide, descrittori, livelli).

\section{Calibrazione: quanto è riuscito l'allineamento dei gain}
Prima di leggere qualunque differenza sulla distanza devo sapere quanto bene sono partiti
allineati i microfoni. Ho misurato l'RMS di ciascuno sul grappolo a 1 m, dove le capsule
erano equidistanti dalla sorgente e dovevano restituire lo stesso livello. La dispersione
è piccola: la deviazione standard fra i microfoni è di 0,17 dB e lo scarto fra il più alto
e il più basso (max-min) è di 0,58 dB. Per canale, lo scarto dalla media va da $-0{,}34$ dB
(ch5) a $+0{,}24$ dB (ch6); gli altri stanno sotto i due decimi di decibel
(figura~\ref{fig:calibrazione}).

La taratura a mano, insomma, è riuscita molto bene. La conseguenza pratica è importante:
ogni differenza per-canale che superi circa 0,6 dB sui segnali successivi è un fatto reale,
non un residuo di calibrazione. Questa è la mia soglia di leggibilità, il fondo sotto il
quale non mi fido di nessuna differenza.
\begin{figure}[ht]
  \centering
  \includegraphics[width=0.7\linewidth]{calibrazione_rms.png}
  \caption{Scarto dell'RMS di ciascun microfono dalla media, sul grappolo a 1 m:
  misura della bonta' dell'allineamento dei gain regolati a mano.}
  \label{fig:calibrazione}
\end{figure}

\section{Il rumore sulla distanza}\label{sec:rumore}
Comincio dal rumore rosa, il segnale più controllato. Due cose mi interessano: come cala
il livello con la distanza, e se il centroide si sposta.

\paragraph{Attenuazione del livello.} Riferito al microfono a 1 m, il livello cala con la
distanza molto meno di quanto vorrebbe il campo libero (figura~\ref{fig:attenuazione}). In
orizzontale misuro $-4{,}6$ dB a 2 m, $-6{,}7$ a 4 m, fino a $-8{,}6$ dB a 8 m; in verticale
i numeri sono più piccoli ($-3{,}3$ a 2 m, $-6{,}7$ a 8 m). Il campo libero di una sorgente
puntiforme darebbe $-6$ dB a ogni raddoppio, cioè $-18$ dB a 8 m, molto più di quanto misuro.
La legge della sala è un'altra, più dolce: circa $-3$ dB a ogni raddoppio, cioè un'ampiezza
che si riduce di $1/\sqrt{2}$ ogni volta che la distanza raddoppia. Mediato sui tre raddoppi,
il mio dato orizzontale fa $-2{,}9$ dB a raddoppio, vicinissimo a quel $-3$.

E qui ho una conferma indipendente. Durante la registrazione ho letto il livello anche al
fonometro (curva A): 96 dBA a 1 m, 93 a 2 m, 90 a 4 m, 87 a 8 m, cioè esattamente $-3$ dB a
ogni raddoppio. Due strumenti diversi, la misura digitale per canale e il fonometro in sala,
danno la stessa legge: il LEAP, sala reale e riverberante, attenua il rumore come
$1/\sqrt{r}$ ($-3$ dB per raddoppio), non come il campo libero $1/r$ ($-6$ dB). Il campo
riverberante riempie lo spazio e ammorbidisce la caduta. Nella figura i punti del fonometro
cadono sulla retta di $-3$ dB e le curve misurate le stanno attorno, ben sopra il campo libero.

\paragraph{Centroide del rumore.} Il centroide del rumore cala nettamente con la distanza
(figura~\ref{fig:centroide-noise}). In orizzontale scende da circa 9000 Hz a 1 m a 5700 Hz
a 8 m, oltre il 30\% in meno; in verticale da 6300 a 5000 Hz. È l'assorbimento dell'aria
sugli acuti: l'aria assorbe le alte frequenze tanto più quanto è lungo il cammino, così i
microfoni lontani ricevono uno spettro più povero di acuti e il loro centroide è più grave.
La doppia ripresa lo conferma: in orizzontale il centroide parte molto più alto che in
verticale (circa 9000 contro 6300 Hz a 1 m), perché la direttività dell'OM1, un fenomeno di
alta frequenza, fa sì che la capsula rivolta verso la sorgente raccolga più acuti.

\paragraph{Il confronto col tetto a 10 kHz.} Questa caduta si vede solo guardando tutta la
banda. Se limito l'analisi a 10 kHz (come in una prima versione, e come nei descrittori
monolitici usati in tempo reale), il centroide del rumore appare quasi piatto, fra 3300 e
3600 Hz, e l'assorbimento dell'aria sparisce (figura~\ref{fig:banda}). Il motivo è che
l'aria agisce soprattutto sopra i 10 kHz: tutto l'effetto della distanza sta nella banda che
il tetto tagliava via. Il tetto a 10 kHz non toglieva segnale inutile, toglieva la firma
della distanza. Sul contrabbasso, invece, le due bande danno quasi lo stesso centroide (lo
strumento ha poca energia sopra i 10 kHz) e l'onda nodo/ventre resta identica: è un fenomeno
della fondamentale, non degli acuti.
\begin{figure}[ht]
  \centering
  \includegraphics[width=0.8\linewidth]{attenuazione_distanza.png}
  \caption{Livello del rumore relativo al microfono a 1 m in funzione della distanza (scala
  log), orizzontale e verticale. I quadrati sono la misura al fonometro in sala; le rette di
  riferimento sono il campo libero ($-6$\,dB per raddoppio) e la legge $1/\sqrt{2}$
  ($-3$\,dB per raddoppio). Misura e fonometro seguono quest'ultima.}
  \label{fig:attenuazione}
\end{figure}
\begin{figure}[ht]
  \centering
  \includegraphics[width=0.8\linewidth]{centroidi_medie_noise.png}
  \caption{Centroide spettrale medio del rumore per canale (distanza), media sul file intero,
  a banda piena: cala con la distanza (assorbimento dell'aria sugli acuti).}
  \label{fig:centroide-noise}
\end{figure}
\begin{figure}[ht]
  \centering
  \includegraphics[width=\linewidth]{centroide_rumore_banda.png}
  \caption{Centroide del rumore in funzione della distanza, a 10 kHz (tratteggio) e a banda
  piena fino a 48 kHz (continuo). Dentro i 10 kHz il centroide è piatto; a banda piena cala
  nettamente con la distanza, l'assorbimento dell'aria sugli acuti che il tetto nascondeva.}
  \label{fig:banda}
\end{figure}

\paragraph{Gli altri descrittori sul rumore.} La distanza non la legge solo il centroide:
quasi tutti i quindici descrittori si muovono con essa, e nello stesso modo in orizzontale e
in verticale (figura~\ref{fig:noise-descrittori}). Il segnale lontano diventa più ``a picchi''
e più tonale. Il crest factor sale forte (da circa 18 a 1 m a circa 32 a 8 m), il tonal power
ratio cresce (da circa $-6$ a $-4$ dB), il numero di picchi passa da 148 a 171, mentre
flatness ed entropy calano. Diventa anche più stazionario: il flux crolla a circa un quarto
(da 0,7 a 0,18). E soprattutto perde gli acuti: zcr e rolloff scendono nettamente (il rolloff
da circa 18 a 11,5 kHz), in linea con il forte calo del centroide. La lettura è coerente:
allontanandosi, l'aria toglie gli acuti a banda larga e resta soprattutto la parte bassa e le
risonanze del luogo, così il rumore lontano è più scuro, più risonante e più fermo. La firma
della distanza è scritta in tutto il vettore, non nel solo centroide.
\begin{figure}[ht]
  \centering
  \includegraphics[width=\linewidth]{medie_noise.png}
  \caption{I sedici descrittori del rumore in funzione della distanza (canale = metri),
  per famiglia, nelle due riprese orizzontale e verticale. Crest, tonal power ratio, numero
  di picchi salgono con la distanza; centroide, rolloff, flatness, entropy, flux, zcr
  scendono.}
  \label{fig:noise-descrittori}
\end{figure}

\section{Il centroide sugli strumenti}
Sugli strumenti il quadro cambia. Sul clarinetto contrabbasso il centroide disegna
un'onda netta lungo i microfoni (figure~\ref{fig:profilo-onda} e~\ref{fig:heatmap}):
non scende in modo regolare con la distanza, ma cresce fino a una cresta al microfono a
3 m (circa $+220$/$+285$ Hz sopra la media del file), poi cala fino a un avvallamento al
microfono a 6 m (circa $-160$/$-190$ Hz), per risalire all'estremo. In numeri assoluti il
centroide del contrabbasso passa da circa 1300 Hz a 1 m a un massimo di circa 1600 Hz a
3 m, scende verso 1180 Hz a 6 m e torna sopra i 1360 Hz a 8 m. Per questi profili uso i soli
cluster e lascio fuori il Do, una nota sola che riprendo fra i gesti.

La cosa più solida è che l'onda è quasi identica in orizzontale e in verticale: ruotando
le capsule la forma non cambia. Questo esclude che sia principalmente un effetto della
direttività del microfono, che cambierebbe con l'orientamento; è qualcosa che sta nella
propagazione e nel luogo. I due gruppi del contrabbasso in orizzontale, registrati con il
trombino originale e con il PETALONIO, si comportano in modo simile fra loro; senza dati
separati non spingo oltre il confronto fra le due campane.

Sul clarinetto soprano l'onda non c'è: il profilo per-canale resta sostanzialmente piatto,
entro circa $\pm 80$ Hz attorno a una media di circa 1740 Hz. Quindi la struttura
per-canale non è una proprietà fissa del luogo o della catena, ma dipende dalla sorgente.
Il soprano ha, come atteso, un centroide più alto del contrabbasso (circa 1740 Hz contro
un valore fra 1180 e 1600 Hz). Quest'onda nasce da un'onda stazionaria nello spazio
acustico (se sia un modo fisso della stanza o una figura legata alla posizione dello
strumento resta da chiarire, e lo riprendo nella discussione).
\begin{figure}[ht]
  \centering
  \includegraphics[width=\linewidth]{profilo_onda_separati.png}
  \caption{Profilo per canale del centroide, deviazione media dalla media del file, su tre
  pannelli a scala verticale condivisa, soli cluster: contrabbasso orizzontale, contrabbasso
  verticale e soprano. Sul contrabbasso la stessa onda nelle due riprese (cresta al microfono
  a 3 m, avvallamento a quello a 6 m); sul soprano resta piatto.}
  \label{fig:profilo-onda}
\end{figure}
\begin{figure}[ht]
  \centering
  \includegraphics[width=0.85\linewidth]{centroidi_heatmap.png}
  \caption{Deviazione del centroide dalla media del file, per canale, su tutti i campioni
  (media sul file intero; i Do sull'intera nota).}
  \label{fig:heatmap}
\end{figure}

\section{Comportamento dei descrittori}
Il centroide è stato il protagonista, ma sedici descrittori non sono sedici doppioni.
Ho guardato come ciascuno si comporta sui fattori dell'esperimento: la distanza, il gesto
(Do grave o cluster), la campana, la sorgente e l'orientamento della capsula. La cosa che
emerge è che ogni descrittore ha un suo profilo di sensibilità: nessuno vede tutto, e si
dividono il lavoro. Comincio da quelli legati al centroide.

\paragraph{Distanza: chi porta l'onda.} Per misurare quanti descrittori condividono l'onda
nodo/ventre del centroide ho correlato, sul contrabbasso, il profilo per canale di ognuno
con quello del centroide (figura~\ref{fig:correlazioni}). Sui cluster la portano in dodici,
centroide compreso. In fase con il centroide (più acuti al nodo) stanno zcr, rolloff,
skewness, entropy, spread, irregularity e tonality, con correlazioni fra 0,92 e 0,98; in
opposizione stanno tpr, flatness, slope e crest (da $-0{,}82$ a $-0{,}96$): dove la
fondamentale è debole questi misurano meno dominanza tonale e meno picco. Restano fuori
kurtosis, flux, obsir-std e numero di picchi (correlazione sotto 0,5 in modulo). La lettura è semplice: il nodo e il ventre sono un'unica leva, la forza della
fondamentale, e tutto ciò che pesa l'equilibrio fra fondamentale e resto si muove insieme.
Per questo l'onda taglia di traverso le quattro famiglie: dentro la forma il centroide la
porta ma obsir-std no, dentro la dinamica zcr e irregularity la portano ma flux no. La domanda
giusta non è a quale famiglia appartiene un descrittore, ma se risponde o no alla fondamentale.

\paragraph{Gesto: dal Do al cluster.} Passare dalla nota singola al cluster muove quasi tutti
i descrittori nello stesso verso. Il cluster è più brillante (centroide da 952 a 1310 Hz sul
PETALONIO, rolloff da 1670 a 2210 Hz, zcr quasi raddoppiato), più denso (numero di picchi da
140 a 151, obsir-std da 0,64 a 1,6, flux più che doppio) e meno dominato da un singolo picco
(crest da 47 a 40, kurtosis da 25 a 13, skewness in calo). È una segnatura coerente del gesto,
leggibile su gran parte del vettore.

\paragraph{Campana: il caso della kurtosis.} Sul Do tenuto la campana lascia un segno reale:
il PETALONIO è più scuro e più appuntito del trombino (centroide del Do 952 contro 1250 Hz,
rolloff 1670 contro 2360 Hz). Il caso più vistoso è la kurtosis (figura~\ref{fig:kurtosis-do}).
Sul Do, in orizzontale, il PETALONIO la fa esplodere al ventre: da 21 al nodo (microfono a
3 m) a 58 e 61 ai microfoni a 5 e 6 m, circa tre volte tanto. Il trombino resta attorno a
10, e il PETALONIO in verticale appiattisce l'effetto (attorno a 13). Serve la congiura di tre
condizioni: la nota singola (sul cluster, con molti picchi, l'effetto svanisce), il ventre
(dove la fondamentale è forte) e la campana verticale puntata di fronte alla capsula. Sui
cluster la kurtosis non segue l'onda nodo/ventre (la sua correlazione col centroide è vicina
a zero): resta indifferente alla distanza. Ma sul Do esplode. Non è dunque un descrittore
della distanza, è un sensore del gesto e della campana: dorme sui cluster e scatta sul Do,
dove diventa la spia più forte dell'interazione fra campana, luogo e microfono.

\paragraph{Sorgente: contrabbasso e soprano.} Il clarinetto soprano si stacca dal contrabbasso
soprattutto in brillantezza e tonalità: centroide attorno a 1740 Hz contro circa 1350,
rolloff 2820 Hz, slope $-2{,}7$ contro $-5$, e soprattutto tpr
18 contro circa 6 e numero di picchi 217 contro 155. I salti grandi che si vedono nelle figure
di famiglia su tpr e numero di picchi sono cambiamenti di sorgente, non di distanza.

\paragraph{Orientamento: dove si vede il microfono.} Ruotando la capsula da orizzontale a
verticale il centroide e i suoi vicini quasi non cambiano (centroide, spread e rolloff entro
il 4\%), coerente con il fatto che l'onda è del luogo e non del trasduttore. In effetti sul
contrabbasso nessun descrittore cambia oltre l'8\% ruotando la capsula: lo strumento ha poca
energia in alto, dove la direttività dell'OM1 agisce. La direttività si fa sentire altrove,
sul rumore a banda larga: lì il centroide orizzontale parte molto più alto del verticale
(§\ref{sec:rumore}), perché la capsula rivolta verso la sorgente raccoglie più acuti.

In conclusione, nessun descrittore vede tutto: il centroide legge distanza e gesto ma è cieco
all'orientamento; flux e obsir-std rispondono soprattutto al gesto e restano fuori dall'onda;
la kurtosis è indifferente all'onda ma esplode sul Do, dove fa da spia della campana. È questa
divisione del lavoro a fare dei sedici un insieme. Per l'uso in tempo reale conviene scegliere
il descrittore per il compito: il centroide per distinguere i gesti restando stabile rispetto
a dove si mette il microfono, l'onda per leggere il luogo, la kurtosis per cogliere il Do e
la campana.

\begin{figure}[ht]
  \centering
  \includegraphics[width=0.9\linewidth]{onda_correlazioni.png}
  \caption{Correlazione del profilo per canale di ciascun descrittore con quello del centroide,
  sul contrabbasso: misura quanto ogni descrittore porta l'onda nodo/ventre. In blu chi la
  segue in fase, in rosso in opposizione, in grigio gli indifferenti (linee tratteggiate a
  $\pm 0{,}8$).}
  \label{fig:correlazioni}
\end{figure}
\begin{figure}[ht]
  \centering
  \includegraphics[width=0.85\linewidth]{kurtosis_do.png}
  \caption{Kurtosis spettrale del Do grave per canale, nelle quattro combinazioni di campana
  e orientamento: il PETALONIO orizzontale esplode al ventre (microfoni a 5 e 6 m), gli altri
  restano piatti.}
  \label{fig:kurtosis-do}
\end{figure}
\begin{figure}[ht]
  \centering
  \includegraphics[width=\linewidth]{medie_forma.png}
  \caption{Famiglia forma: medie per canale dei descrittori, per campione.}
  \label{fig:forma}
\end{figure}
\begin{figure}[ht]
  \centering
  \includegraphics[width=\linewidth]{medie_distribuzione.png}
  \caption{Famiglia distribuzione: medie per canale dei descrittori, per campione.}
  \label{fig:distribuzione}
\end{figure}
\begin{figure}[ht]
  \centering
  \includegraphics[width=\linewidth]{medie_tonalita.png}
  \caption{Famiglia tonalita': medie per canale dei descrittori, per campione.}
  \label{fig:tonalita}
\end{figure}
\begin{figure}[ht]
  \centering
  \includegraphics[width=\linewidth]{medie_dinamica.png}
  \caption{Famiglia dinamica: medie per canale dei descrittori, per campione.}
  \label{fig:dinamica}
\end{figure}

\section{Discussione}
Il filo dell'esperimento è una discrepanza oggettiva: lo stesso evento sonoro, ripreso da
punti diversi dello spazio acustico, restituisce descrittori diversi, e la differenza
supera nettamente la soglia di leggibilità fissata in calibrazione. Sul rumore il centroide
cala con la distanza (l'aria assorbe gli acuti); sul contrabbasso disegna un'onda di centinaia
di Hz da un microfono all'altro. In nessuno dei due casi è un effetto di livello, perché il
centroide è invariante di scala: è la forma dello spettro che arriva diversa a ogni microfono.

La doppia ripresa, orizzontale e verticale, è la chiave per leggere questa differenza. Se
l'onda fosse un artefatto del microfono (la direttività dell'OM1 ai medio-acuti) cambierebbe
ruotando le capsule; invece resta quasi identica nei due orientamenti. Questo mi porta a
considerare sorgente, ambiente e catena come un sistema unico: il valore che leggo non è una
proprietà della sola sorgente, ma il risultato di come quella sorgente si propaga in quel
luogo e viene raccolta in quel punto. Il fatto che l'onda compaia sul contrabbasso e non sul
soprano dice che anche la sorgente entra nel conto: non tutto si spiega con il luogo.

L'onda del centroide sul contrabbasso ha un'origine che ho cominciato a leggere già in sala.
Lungo i microfoni la fondamentale del clarinetto contrabbasso non arriva con la stessa forza:
è debole dalle parti del microfono a 3 m e forte attorno ai microfoni a 5 e 6 m, il segno di
un'onda stazionaria nello spazio acustico. Dove la fondamentale ha un nodo è debole, e il
centroide, che pesa le frequenze per la loro ampiezza, si sposta verso gli acuti; dove ha un
ventre è forte e tira il centroide verso il basso. È esattamente ciò che misuro: la cresta
del centroide al microfono a 3 m e l'avvallamento a quello a 6 m, identici in orizzontale e
in verticale, quindi non è il microfono.

Resta però un dubbio che non chiudo con questi dati. Non so ancora se quel nodo sia un modo
proprio della stanza, fisso nello spazio a prescindere dalla sorgente, oppure una figura
d'interferenza legata alla posizione del clarinetto (il suono diretto più una riflessione
forte), che si sposterebbe insieme alla sorgente. La prova pulita è una sola: spostare la
sorgente e rifare la misura. Se il nodo torna nello stesso punto è un modo della stanza; se
si sposta con lo strumento è legato alla geometria fra sorgente e microfoni. È il controllo
che qui manca, e il prossimo passo naturale dell'esperimento. Il soprano, intanto, non
disegna la stessa onda: la sua fondamentale più acuta dispone nodi e ventri a un passo più
fitto e in posizioni diverse, che sul cluster non si sommano in una figura altrettanto netta.

\section{Conclusioni}
Ho imparato che i descrittori non sono indifferenti al punto di ripresa. La calibrazione a
mano è ottima (sotto 0,6 dB di scarto fra i microfoni), quindi le differenze che ho letto
sono reali. Il livello cala con la distanza ma molto meno del campo libero, perché il LEAP è
una sala riverberante. A banda piena il centroide del rumore cala nettamente con la distanza,
l'assorbimento dell'aria sugli acuti che il vecchio tetto a 10 kHz nascondeva, e con lui si
muovono gli altri descrittori del rumore (il segnale lontano è più scuro, più tonale e più
stazionario). Sul contrabbasso,
però, il centroide e con lui i descrittori di forma (e parte della dinamica) disegnano
un'onda netta lungo i microfoni, uguale nelle due direzioni di ripresa: una deformazione
dello spettro che dipende dalla propagazione e dal luogo, non dalla direttività del microfono.
La tonalità e parte della distribuzione restano più ferme.

Per l'uso in tempo reale, nel sistema di \emph{interfantasia}, la conseguenza è concreta. Se
i descrittori che guidano il comportamento dipendono anche dal punto in cui si raccoglie il
segnale, allora la posizione del microfono fa parte della calibrazione tanto quanto il
guadagno: spostarlo significa spostare il vettore. I descrittori di forma, centroide in testa,
sono i più sensibili al luogo e alla catena, e vanno trattati con questa consapevolezza;
quelli più fermi (tonalità) offrono un appoggio più stabile. La firma del luogo, insomma,
entra nella misura: conviene conoscerla prima di affidarle il controllo.

\appendix
\section{Note matematiche}
\subsection{Il centroide come invariante di scala}
Il centroide spettrale è il baricentro dello spettro di ampiezza:
\begin{equation}
  C = \frac{\sum_k f_k\, |X_k|}{\sum_k |X_k|}.
\end{equation}
Moltiplicare lo spettro per una costante $\alpha>0$ (un guadagno) lascia $C$ invariato:
\begin{equation}
  \frac{\sum_k f_k\, \alpha|X_k|}{\sum_k \alpha|X_k|}
  = \frac{\alpha \sum_k f_k\, |X_k|}{\alpha \sum_k |X_k|} = C.
\end{equation}
Quindi una variazione del centroide tra due microfoni non può essere un effetto di
livello: segnala una deformazione della \emph{forma} dello spettro.

\subsection{Legge dell'inverso della distanza}
In campo libero la pressione efficace di una sorgente puntiforme decresce come $1/r$.
In decibel, al raddoppio della distanza:
\begin{equation}
  \Delta L = 20\log_{10}\!\frac{1/(2r)}{1/r} = 20\log_{10}\tfrac{1}{2} \approx -6{,}02\ \text{dB}.
\end{equation}
La sala reale del LEAP non è anecoica, e infatti non segue questa retta. La misura, sia per
canale sia al fonometro, cade su una legge più dolce, di $-3$ dB a ogni raddoppio:
\begin{equation}
  \Delta L = 20\log_{10}\frac{1}{\sqrt{2}} \approx -3{,}01\ \text{dB},
\end{equation}
cioè un'ampiezza proporzionale a $1/\sqrt{r}$. È il comportamento di un campo non più
puramente diretto: il contributo riverberante della stanza, che decade poco con la distanza,
si somma al diretto e ne ammorbidisce la caduta.

\subsection{Assorbimento dell'aria sugli acuti}
L'assorbimento atmosferico cresce con la frequenza e con il cammino percorso: a parità di
sorgente, i microfoni più lontani ricevono uno spettro relativamente più povero di acuti.
Poiché il centroide pesa le frequenze per la loro ampiezza, una perdita selettiva degli
acuti lo sposta verso il basso al crescere della distanza (effetto qualitativo, non un
adattamento quantitativo).

\end{document}
